\documentclass[11pt]{article}
\usepackage[margin=1in]{geometry}
\usepackage{amsmath}
\usepackage{amssymb}

\setlength{\parindent}{0pt}
\setlength{\parskip}{0.6em}

\begin{document}

\section*{Review of \texttt{main.tex}: SENA manuscript}

\textbf{Overall assessment.}
The manuscript is ambitious, substantially more implementation-aware than a typical concept note, and mostly internally coherent in its core inventory transition and LogNormal lead-time mapping.
I explicitly checked symbol definitions, sign conventions, dimensional consistency, and implementation-facing formulas.
I did not find inverse-trig or quadrant/branch issues because the manuscript does not use that class of expression, and I did not find a clear sign/factor error in the core inventory or lead-time equations.
The highest-priority problems are instead a small number of concrete specification gaps and one true state-space bug that would change the model if implemented literally.

\section*{Most important technical findings}

\begin{enumerate}
  \item \textbf{Definite error: a positive initial pipeline $B^{\mathrm{init}}_i$ cannot ever generate receipts.}

  \textbf{Location:} the initial-state block in ``What SENA models'' and the receipt-assignment formula in ``Restocking as a policy with an explicit order pipeline.''

  \textbf{Problem:} the manuscript allows
  \[
    B_{1,i}=B^{\mathrm{init}}_i,\qquad B^{\mathrm{init}}_i\ge 0,
  \]
  but later defines receipts only through post-$t_1$ order batches,
  \[
    R_{k,i}=\sum_{u=2}^{k} O_{u,i}\,\mathbf{1}\{t_{k-1}<a_{u,i}\le t_k\}.
  \]
  There is no latent pre-$t_1$ order history or residual-arrival distribution tied to $B^{\mathrm{init}}_i$.
  If $B^{\mathrm{init}}_i>0$, those units have no mechanism to arrive.

  \textbf{Why it matters:} early-interval receipts, pipeline depth, reorder propensity, and stockout risk are all biased if the model literally permits a positive initial in-transit state that is never receivable.
  This is not just a missing prior; it is a state-transition inconsistency.

  \textbf{Suggested fix:} either make $B^{\mathrm{init}}_i=0$ an enforced model assumption, or introduce latent pre-$t_1$ pseudo-orders / a residual arrival-time distribution whose outstanding mass equals $B^{\mathrm{init}}_i$.

  \item \textbf{Definite error: the order-size location parameter $\mu^{O}_{k,i}$ is used but never defined.}

  \textbf{Location:} the LogNormal order-size model in the replenishment section and the later ``Order size model'' discussion in the hyperparameter block.

  \textbf{Problem:} the manuscript specifies
  \[
    O_{k,i}\mid(G_{k,i}=1)\sim\text{LogNormal}(\mu^{O}_{k,i},\sigma^{O}_i),
  \]
  but never defines what determines $\mu^{O}_{k,i}$, whether it is static or dynamic, whether it depends on the posterior gap to a target stock level, or what prior/transition it has.

  \textbf{Why it matters:} this makes the generative model incomplete and leaves order-size inference non-reproducible.
  Two implementations could choose fundamentally different order-quantity dynamics while both claiming to match the manuscript.

  \textbf{Suggested fix:} define $\mu^{O}_{k,i}$ explicitly.
  A natural option is to link it to a target inventory-position gap, e.g.
  \[
    \mu^{O}_{k,i}=\alpha^{O}_i+\beta^{O}_i\,\bigl(T_{k,i}-x_{k-1,i}-B_{k-1,i}\bigr)_+,
  \]
  or to a simpler SKU-specific baseline with clear priors if the paper wants a less policy-driven replenishment model.

  \item \textbf{Definite mismatch: the lead-time variability observation is connected to the wrong latent quantity.}

  \textbf{Location:} the ``Price and lead-time observations'' block, specifically the relative-width definition
  \[
    w_{k,i}=\frac{\tilde L^{\mathrm{hi}}_{k,i}-\tilde L^{\mathrm{lo}}_{k,i}}{\max\!\left\{\frac{\tilde L^{\mathrm{hi}}_{k,i}+\tilde L^{\mathrm{lo}}_{k,i}}{2},\ \epsilon_L\right\}},
  \]
  followed by the ordered-logit model on $\psi_{k,i}=\log v^L_{k,i}$.

  \textbf{Problem:} $w_{k,i}$ is a \emph{scale-free relative-width} summary, but the observation model uses raw log-variance $\psi_{k,i}=\log v^L_{k,i}$ as its latent driver:
  \[
    \tilde c^L_{k,i}\sim\text{OrderedLogit}\!\bigl(a_{0,i}+a_{1,i}\psi_{k^{c,\star}_{k,i},i},\ \kappa_{1:C-1}\bigr).
  \]
  That means two SKUs with the same relative lead-time uncertainty but different mean lead times would imply the same observed range class and different latent predictors.

  \textbf{Why it matters:} the model will systematically confound mean lead time with variability class.
  This directly affects posterior uncertainty in arrivals and therefore stockout-risk and reorder calculations.

  \textbf{Suggested fix:} drive the ordinal variability class with a scale-free quantity such as the LogNormal shape parameter $\sigma_{L,k,i}$ from the mean/variance map, or equivalently with
  \[
    \log\!\left(\frac{v^L_{k,i}}{(m^L_{k,i})^2}\right)
  \]
  (or $\log(1+v^L_{k,i}/(m^L_{k,i})^2)$ if you want exact compatibility with the LogNormal parameterization).

  \item \textbf{Likely issue: the reorder trigger is not pipeline-aware, even though the later reorder-quantity rule is.}

  \textbf{Location:} the planning outputs block, especially the trigger
  \[
    P\bigl(x_{k,i}\le \text{ROP}_{k,i}\mid\text{data}\bigr)
  \]
  versus the later inventory-position definition
  \[
    IP_{k,i}=x_{k,i}+B_{k,i}
  \]
  and need probability
  \[
    P\!\left(D^H_{k,i}>IP_{k,i}\mid\text{data}\right).
  \]

  \textbf{Problem:} the manuscript first defines a reorder trigger on \emph{on-hand} inventory only, but later defines reorder quantity from \emph{inventory position} (on-hand plus in-transit pipeline).
  Those can disagree materially when there is already stock on order.

  \textbf{Why it matters:} a low on-hand state with a large pipeline can trigger a reorder under the first rule even when the second rule correctly says no new order is needed.
  That inconsistency is especially important because SENA explicitly models the pipeline rather than treating replenishment as instantaneous.

  \textbf{Suggested fix:} either replace the trigger with
  \[
    P\bigl(IP_{k,i}\le \text{ROP}_{k,i}\mid\text{data}\bigr),
  \]
  or state clearly that the on-hand trigger is only a simplified no-open-order heuristic and that the inventory-position rule is the recommended operational output.

  \item \textbf{Likely issue: within-interval stockout quantities are not fully defined by the current generative model.}

  \textbf{Location:} the stockout-flag observation block and the output definition
  \[
    \text{SOR}_{k,i}\equiv P\Bigl(\min_{t\in[t_{k-1},t_k]}x_i(t)=0\ \Big|\ \text{data}\Bigr).
  \]

  \textbf{Problem:} receipt times have an explicit within-interval timestamp through $a_{u,i}$, but demand $D_{k,i}$ and corrections $A^{\mathrm{adj}}_{k,i}$ are modeled as interval aggregates.
  The manuscript never specifies a within-interval arrival process for demand or a timing convention for adjustments.
  As written, quantities involving $\min_t x_i(t)$ and stockout-flag likelihoods depend on unstated path assumptions.

  \textbf{Why it matters:} stockout-risk outputs and the interpretation of stockout flags become implementation-dependent.
  Two teams could use the same interval-level posterior draws and obtain different $\text{SOR}_{k,i}$ values just by choosing different intra-interval interpolation rules.

  \textbf{Suggested fix:} specify one explicit within-interval approximation, e.g.
  an inhomogeneous Poisson or renewal demand-arrival process inside $\Delta t_k$, receipt jumps at $a_{u,i}$, and a clear timing convention for $A^{\mathrm{adj}}_{k,i}$.
  If that is outside scope, relabel $\text{SOR}_{k,i}$ as a proxy rather than a fully defined posterior quantity.

  \item \textbf{Likely issue: the abrupt-drift block is still a modeling intention, not yet a generative specification.}

  \textbf{Location:} ``Abrupt drift via occasional change points.''

  \textbf{Problem:} the paper introduces
  \[
    c_k\sim\text{Bernoulli}(p_c)
  \]
  and says that $c_k=1$ allows a partial reset or jump, but it does not write the actual transition equations for the affected states, the reset distribution, or which components are permitted to jump.

  \textbf{Why it matters:} this prevents reproducible implementation and makes the change-point ablation under-specified.
  In the current form, the manuscript describes the intention to model abrupt drift without yet fully defining the model.

  \textbf{Suggested fix:} add one compact mixture-transition block, for example for $\log\lambda_{k,j}$ (and any other jumping state), with explicit distributions under $c_k=0$ and $c_k=1$.
\end{enumerate}

\section*{Notation drift snapshot (required check)}

\begin{itemize}
  \item \textbf{$\beta$.} The symbol first appears as the time-varying stock-sensitivity coefficient $\beta_{k,i}$ in the order-placement policy, later as the price elasticities $\beta^{\mathrm{svc}}_j$ and $\beta^{\mathrm{ret}}_i$, and later again as the target cycle service level in the planning section.
  Even though the equations are locally readable, the same symbol is carrying three unrelated meanings.
  \textbf{Suggested fix:} rename the service-level target to something like $\alpha_{\mathrm{SL}}$ or $p_{\mathrm{SL}}$.

  \item \textbf{$\tau$.} The symbol is used first as the ranking-temperature parameter in the ordinal likelihood and later as the gating threshold in the reorder recommendation.
  \textbf{Suggested fix:} keep $\tau$ (or $\tau_{\mathrm{rank}}$) for ranking noise and rename the reorder gate to $\tau_{\mathrm{gate}}$ or $p_0$.
\end{itemize}

\section*{Meaningful editorial findings}

\begin{enumerate}
  \item \textbf{Location:} abstract and contribution-summary language versus the later ``Evaluation protocol for publication readiness'' section.

  \textbf{Problem:} the opening pages read like a validated method description, but the manuscript later makes clear that empirical validation is still proposed rather than reported.

  \textbf{Why it matters:} this creates avoidable claim-strength tension.
  Readers may reasonably expect experimental evidence that is not actually present in the draft.

  \textbf{Suggested fix:} mark the paper more explicitly as a proposed model specification until the evaluation section is executed.

  \textbf{Candidate rewrite:} ``SENA is a proposed Bayesian generative model for sparse-stock settings; its intended outputs are posterior inventory trajectories, demand-rate estimates, stockout-risk summaries, and replenishment recommendations under uncertainty.''

  \item \textbf{Location:} opening of ``Inventory dynamics.''

  \textbf{Problem:} the setup around the first transition equations is awkwardly split across two short framing sentences:
  ``For each SKU $i$, with initial states $(x_{1,i},B_{1,i})$ and transitions over $k=2,\dots,K$:'' immediately followed by ``For all recursions in this section, $k=2,\dots,K$ unless otherwise stated.''

  \textbf{Why it matters:} this is small, but it makes the key model-definition section feel less polished than the mathematics deserves.

  \textbf{Suggested fix:} collapse the frame into one sentence.

  \textbf{Candidate rewrite:} ``For each SKU $i$, with initial states $(x_{1,i},B_{1,i})$, the transition recursions below apply for $k=2,\dots,K$ unless otherwise stated.''

  \item \textbf{Location:} adjustment hyperparameter description in the default-priors section.

  \textbf{Problem:} the text describes $s^{\mathrm{adj}}_{r,i}$ as being ``in SKU units,'' but the model uses
  \[
    A^{\mathrm{adj}}_{k,i}\sim\text{Normal}\bigl(m^{\mathrm{adj}}_{r,i}\,\Delta t_k,(s^{\mathrm{adj}}_{r,i})^2\,\Delta t_k\bigr),
  \]
  so the implied standard deviation is $s^{\mathrm{adj}}_{r,i}\sqrt{\Delta t_k}$.

  \textbf{Why it matters:} this affects how readers elicit priors and interpret units under irregular intervals.

  \textbf{Suggested fix:} describe $s^{\mathrm{adj}}_{r,i}$ as a diffusion-scale parameter, or state explicitly that the interval-level standard deviation is $s^{\mathrm{adj}}_{r,i}\sqrt{\Delta t_k}$.

  \textbf{Candidate rewrite:} ``Here $s^{\mathrm{adj}}_{r,i}>0$ is the unit-time diffusion scale for corrections; over interval $\Delta t_k$, the correction standard deviation is $s^{\mathrm{adj}}_{r,i}\sqrt{\Delta t_k}$.''
\end{enumerate}

\section*{Proofing sweep (required scan)}

I ran the required scan on \texttt{main.pdf}, \texttt{main.tex}, and \texttt{refs.bib}, then spot-checked the hits.
The PDF scan was dominated by false positives from table-of-contents dot leaders, so the source-text scan was more informative.
I did not find duplicated-punctuation or malformed-title defects in \texttt{refs.bib}.

\begin{itemize}
  \item \textbf{Location:} concluding sentence of ``Promotion-bundles as explicit services.''

  \textbf{Problem:} the scanner correctly flags the semicolon in ``expected to improve identifiability and interpretability; the evaluation section specifies ablations \dots'' as a semicolon-capitalization seam.

  \textbf{Suggested fix:} split this into two sentences or replace the semicolon with an em dash.

  \item \textbf{Location:} the paragraph immediately following the relative-range equation in the lead-time observation block.

  \textbf{Problem:} the lower-case ``where $\epsilon_L>0$ \dots'' reads like a display-fragment rather than a complete sentence.

  \textbf{Suggested fix:} rewrite as ``Here, $\epsilon_L>0$ is a fixed lead-time floor \dots'' or attach the clause more tightly to the preceding sentence.
\end{itemize}

\section*{Highest-priority fixes}

\begin{enumerate}
  \item Give the initial pipeline a real receipt mechanism, or enforce $B^{\mathrm{init}}_i=0$ throughout the model.
  \item Define $\mu^{O}_{k,i}$ explicitly.
  \item Re-map ordinal lead-time variability inputs to a scale-free latent variability quantity.
  \item Make the reorder trigger pipeline-aware.
  \item Specify the within-interval path assumptions used for stockout risk and stockout flags.
  \item Add explicit post-change transition equations for abrupt drift.
\end{enumerate}

\section*{Items needing external verification}

I did not identify an outside factual claim that is clearly wrong on its face.
The main issues are internal consistency and implementation specificity.
The one class of claims that still needs verification is empirical: statements that promo/bundle modeling improves identifiability or decision quality should stay framed as hypotheses until the evaluation protocol is actually run.

\end{document}